\chapter{Einleitung}
\label{sec: Einleitung}

Im Zusammenhang mit dem Titel der Themenstellung ist zunächst das Ziel der Arbeit zu definieren. Dazu können/sollten vorab Ausgangssituation, Umfeld, bereits bekannte Lösungskonzepte, Schwachstellen und Defizite bestehender technischer und/oder wirtschaftlicher Lösungen, Forderungen des Marktes, Trends, gesteigerte Qualitätsansprüche, Umweltforderungen \usw durchaus als Motivation für das gewählte Thema schlüssig dargestellt und daraus die Aufgabenstellung der Arbeit abgeleitet werden. Auf allfällig für die Bearbeitung der Aufgabenstellung verwendete Methoden kann hingewiesen werden, um den Bezug zu den nachfolgend diskutierten Grundlagen herzustellen.
\begin{itemize}
	\item Umfang: \mx 3 Seiten, nicht zu detailliert
	\item Ziel: Interesse wecken
	\item Schwerpunkt: Ausgangssituation, Umfeld, Problemstellung (sehr ausführlich), Rahmenbedingungen, Schwachstellen bestehender Lösungen
	\item Motivation: Warum gerade diese Arbeit?
	\item Herausforderung: Was ist neu? (Forschungsfrage)
	\item Aufgabenstellung: Wie sollte das Problem gelöst werden?
	\item Ziele und Nichtziele
	\item Lösungsschritte mit Hinweis auf „erreicht“ oder „nicht erreicht“
\end{itemize}